% =============================================================================
%  Tài liệu 04 - Chất lượng và kiểm thử (Quality & Test Document)
%  Build: latexmk 04-quality-and-test.tex   (pdfLaTeX, output in build/)
% =============================================================================
\documentclass[12pt,a4paper]{report}
\usepackage{sumtalk-report}
% =============================================================================
%  sumtalk-info.tex - information shared by every document (01-06).
%  Edit here once; every cover and footer picks it up.
% =============================================================================
\newcommand{\ProjectName}{SumTalk}
\newcommand{\CourseName}{QUẢN LÝ DỰ ÁN PHẦN MỀM}          % cover (upper case)
\newcommand{\CourseNameNormal}{Quản lý Dự án Phần mềm}     % footer
\newcommand{\AcademicYear}{2026 -- 2027}
\newcommand{\Instructor}{\canDien[họ tên giảng viên]}

% One \item per member: "Full name -- student ID", sorted by given name.
\newcommand{\TeamMembers}{%
  \item Nguyễn Anh Duy -- \canDien[MSSV]
  \item Đỗ Thành Đạt -- \canDien[MSSV]
  \item Trần Hải Đăng -- \canDien[MSSV]
  \item Châu Hùng Hào -- \canDien[MSSV]
  \item Bùi Đoàn Minh Nhật -- \canDien[MSSV]
  \item Võ Hoàng Sơn -- \canDien[MSSV]
}


\newcommand{\DocTitle}{CHẤT LƯỢNG VÀ KIỂM THỬ}
\newcommand{\DocTitleNormal}{Chất lượng và kiểm thử}
\newcommand{\DocTitleEn}{Quality \& Test Document}
\newcommand{\DocRunningTitle}{Chất lượng và kiểm thử}
\newcommand{\DocCode}{SUMTALK-QTD}
\newcommand{\DocVersion}{2.1 -- 24/11/2026}
\newcommand{\DocStatus}{Chờ duyệt}
\newcommand{\DocAuthor}{\canDien[Project Manager]}
\newcommand{\DocReviewers}{QC, Product Owner, Project Manager}
\newcommand{\DocReaders}{QC, Product Owner, Project Manager}
\newcommand{\CoverDate}{TP. HỒ CHÍ MINH, THÁNG 11/2026}
\hypersetup{pdftitle={SumTalk -- Chất lượng và kiểm thử},pdfauthor={Nhóm SumTalk}}

% A user-visible label of the app (kept in English, bold, as in the md version).
\newcommand{\ui}[1]{\textbf{#1}}
% "Passed" mark used in the result tables.
\newcommand{\pass}{$\checkmark$}
% UI symbols quoted from the app (pdfLaTeX has no glyph for them by default).
\DeclareUnicodeCharacter{25BE}{\ensuremath{\blacktriangledown}}
\DeclareUnicodeCharacter{2026}{\ldots}

\begin{document}

\makereportcover

% ============================================================ Front matter
\chapter*{Thông tin tài liệu}
\addcontentsline{toc}{chapter}{Thông tin tài liệu}

\docinfotable

\section*{Lịch sử phiên bản}

\begin{table}[H]
  \centering
  \begin{tabularx}{\linewidth}{@{}Q{1.6cm}Q{2.3cm}LP{3.2cm}@{}}
    \toprule
    \thead{Phiên bản} & \thead{Ngày} & \thead{Nội dung thay đổi} & \thead{Người thực hiện} \\
    \midrule
    1.0 & 12/11/2026 & Bản đầu, gồm hai tài liệu riêng: Quality \& Test Plan và Test Summary Report & \canDien \\
    2.0 & 13/11/2026 & Gộp thành một tài liệu, đánh lại số mục (CR-08) & \canDien \\
    2.1 & 24/11/2026 & Chuyển sang định dạng báo cáo PDF (LaTeX), sắp xếp lại thành năm chương và một phụ lục;
      thêm sơ đồ các tầng kiểm thử và biểu đồ kết quả & \canDien \\
    \bottomrule
  \end{tabularx}
\end{table}

\needspace{8\baselineskip}
\heading{Quy ước trong tài liệu}
\begin{itemize}
  \item \canDien[\ldots]: số liệu đã có ở Azure Boards, dashboard chi phí hoặc biên bản họp nhưng chưa chép
        vào tài liệu, hoặc kết quả kiểm thử QC chưa điền; cần bổ sung trước khi trình duyệt.
  \item \textbf{Số liệu kiểm thử} (số test, kết quả đo hiệu năng) lấy từ repo ngày 12/11/2026 trên bản v0.4.
  \item \textbf{Mã story}: \#60, \#67, \#99, \#100 là ID trên Azure Boards; S-xx và T-xx là mã test case của QC
        (Phụ lục~\ref{app:testcases}).
  \item \textbf{Nhãn giao diện} được giữ nguyên tiếng Anh và in đậm (ví dụ \ui{Summary}); Phụ lục~\ref{app:testcases}
        giữ nguyên tiếng Anh như bản QC đang dùng.
\end{itemize}

\tableofcontents

\chapter*{Tóm tắt}
\addcontentsline{toc}{chapter}{Tóm tắt}

Tài liệu này đặt \textbf{kế hoạch chất lượng và kiểm thử} của SumTalk cạnh \textbf{kết quả kiểm thử} để đối
chiếu trực tiếp. Phần kế hoạch xác định năm mục tiêu chất lượng, Definition of Done gồm tám điều kiện cho mọi
story, chiến lược kiểm thử ở tám cấp (từ unit test chạy mỗi commit tới UAT trước phát hành), thang bốn mức lỗi
và các tiêu chí bắt đầu, kết thúc, phát hành.

Phần kết quả trên bản v0.4 (ngày 12/11/2026) cho thấy mọi bộ test tự động đã chạy đều đạt: \textbf{88/88} test
backend, \textbf{8/8} smoke test desktop cả từ mã nguồn lẫn bản đóng gói, \textbf{8/8} kiểm tra Docker, lint
sạch. Hiệu năng đạt mục tiêu với khoảng cách lớn: chữ hiện khoảng \textbf{1\,s} sau khi nói, file 11,6 phút xử
lý xong trong \textbf{17,8--56,7\,s}, app desktop khởi động trong 1,4--1,9\,s. Kết luận: bản v0.4 \textbf{đủ
điều kiện cho đợt test nội bộ và UAT}, nhưng \textbf{chưa đủ để nghiệm thu cuối} vì còn ba việc: chạy lại test
E2E trình duyệt sau thay đổi của Sprint 4, deploy thật lên AWS, và đo độ chính xác tiếng Việt trên mẫu UAT.
Phụ lục liệt kê 23 test case E2E của Sprint 3 cho tóm tắt (\#99) và dịch (\#100).

\textbf{Từ khoá:} quality plan, test plan, Definition of Done, kiểm thử tự động, Playwright, pytest, test summary
report, tiêu chí phát hành, test case E2E.

\chapter*{Danh mục từ viết tắt}
\addcontentsline{toc}{chapter}{Danh mục từ viết tắt}

\begin{center}
  \setstretch{1.1}
  \begin{tabular}{@{}P{2.2cm}P{13.2cm}@{}}
    \toprule
    \thead{Viết tắt} & \thead{Ý nghĩa} \\
    \midrule
    AC     & Acceptance Criteria -- tiêu chí chấp nhận của story \\
    API    & Application Programming Interface -- giao diện lập trình ứng dụng \\
    AWS    & Amazon Web Services (EC2: máy chủ ảo; S3: lưu trữ đối tượng) \\
    BA     & Business Analyst -- chuyên viên phân tích nghiệp vụ \\
    BLEU   & Bilingual Evaluation Understudy -- chỉ số đánh giá bản dịch máy \\
    chrF++ & Character n-gram F-score -- chỉ số đánh giá bản dịch theo n-gram ký tự và từ \\
    CPU    & Central Processing Unit -- bộ xử lý trung tâm \\
    DoD    & Definition of Done -- tiêu chí hoàn thành \\
    E2E    & End-to-End -- kiểm thử đầu cuối \\
    GPU    & Graphics Processing Unit -- bộ xử lý đồ hoạ, dùng để chạy model \\
    HTTPS  & Hypertext Transfer Protocol Secure -- giao thức web có mã hoá \\
    LLM    & Large Language Model -- mô hình ngôn ngữ lớn \\
    NFR    & Non-Functional Requirement -- yêu cầu phi chức năng \\
    PM     & Project Manager -- quản lý dự án \\
    PO     & Product Owner -- chủ sản phẩm \\
    QC     & Quality Control -- kiểm soát chất lượng \\
    ROUGE  & Recall-Oriented Understudy for Gisting Evaluation -- chỉ số đánh giá bản tóm tắt \\
    SRT    & SubRip Text -- định dạng file phụ đề \\
    STT    & Speech-to-Text -- chuyển giọng nói thành văn bản \\
    UAT    & User Acceptance Testing -- kiểm thử chấp nhận của người dùng \\
    WER    & Word Error Rate -- tỷ lệ lỗi từ \\
    \bottomrule
  \end{tabular}
\end{center}

\cleardoublepage\phantomsection
\addcontentsline{toc}{chapter}{\listfigurename}
\listoffigures

\cleardoublepage\phantomsection
\addcontentsline{toc}{chapter}{\listtablename}
\listoftables

% ============================================================ Chapter 1
\chapter{Giới thiệu}
\label{chap:intro}

\section{Bối cảnh}

Đặc tả yêu cầu phần mềm~\cite{srs} đã đặt ra acceptance criteria cho từng user story và 37 yêu cầu phi chức năng
đo được. Kế hoạch quản lý dự án~\cite{pmp} dành một nhánh WBS cho kiểm thử và phát hành, gồm test tự động, test
E2E của Sprint 3 (\#99, \#100) và UAT. Tài liệu này cho biết các yêu cầu đó được kiểm tra bằng cách nào, khi nào
một story được coi là xong, và bản phát hành hiện tại đã đạt tới đâu.

\section{Mục đích và người đọc của tài liệu}

Tài liệu đặt kế hoạch chất lượng và kiểm thử cạnh kết quả kiểm thử để đối chiếu trực tiếp: mục tiêu chất lượng,
Definition of Done, chiến lược kiểm thử, quản lý lỗi, tiêu chí phát hành, rồi kết quả đo được trên bản v0.4 và bộ
test case E2E. Người đọc chính là QC (lập và chạy test, báo lỗi), Product Owner (chấp nhận story theo Definition of
Done, đồng ý phát hành) và Project Manager (theo dõi chất lượng như một phần của tiến độ). Kết luận ở
Chương~\ref{chap:summary} là đầu vào cho quyết định nghiệm thu trong \docref{06}{Báo cáo kết thúc dự
án}~\cite{pcr}.

\section{Cấu trúc tài liệu}

Phần còn lại của tài liệu gồm bốn chương và một phụ lục:
\begin{itemize}
  \item \textbf{Chương~\ref{chap:quality} -- Kế hoạch chất lượng:} mục tiêu chất lượng, Definition of Done và các
        hoạt động rà soát.
  \item \textbf{Chương~\ref{chap:testplan} -- Kế hoạch kiểm thử:} chiến lược kiểm thử theo cấp, quản lý lỗi, tiêu
        chí bắt đầu, kết thúc và phát hành.
  \item \textbf{Chương~\ref{chap:summary} -- Báo cáo tổng kết kiểm thử:} môi trường, kết quả test tự động, hiệu
        năng, test thủ công, lỗi còn tồn và đánh giá theo tiêu chí phát hành.
  \item \textbf{Chương~\ref{chap:conclusion} -- Kết luận và phê duyệt.}
  \item \textbf{Phụ lục~\ref{app:testcases} -- Test case E2E Sprint 3} cho tóm tắt (\#99) và dịch (\#100).
\end{itemize}

% ============================================================ Chapter 2
\chapter{Kế hoạch chất lượng}
\label{chap:quality}

\textbf{Chất lượng được giữ bằng bốn lớp: Definition of Done cho mọi story, test tự động ở ba tầng (pytest,
Playwright trình duyệt, Playwright desktop), test thủ công theo test case của QC, và tiêu chí phát hành rõ
ràng.} Tính đến 12/11/2026, 88 test backend đều đạt, ruff sạch và smoke test desktop đạt 8/8.

\section{Mục tiêu chất lượng}

Năm thuộc tính chất lượng được theo dõi, mỗi thuộc tính có mục tiêu và cách đo cụ thể
(Bảng~\ref{tab:quality-goals}). Mục tiêu về độ chính xác AI và hiệu năng lấy trực tiếp từ các yêu cầu phi chức
năng trong~\cite{srs}.

\begin{table}[htbp]
  \centering
  \caption{Mục tiêu chất lượng}
  \label{tab:quality-goals}
  \small
  \begin{tabularx}{\linewidth}{@{}P{2.9cm}LP{5.4cm}@{}}
    \toprule
    \thead{Thuộc tính} & \thead{Mục tiêu} & \thead{Đo bằng} \\
    \midrule
    Đúng chức năng & 100\% acceptance criteria của các story trong sprint đạt & Test case QC, test E2E \\
    Độ chính xác AI & Theo NFR-A1--A5 trong đặc tả yêu cầu~\cite{srs}
      & \path{eval_summary.py}, \path{eval_translation.py}, WER trên mẫu UAT \\
    Hiệu năng & Live $\leq$ 3--5\,s; 1 giờ audio $\leq$ 10 phút trên GPU & Script đo độ trễ, upload file dài \\
    Ổn định & 0 lỗi Critical/High còn mở khi phát hành & Azure Boards (Bug) \\
    Bảo trì & Lint sạch, test tự động chạy $<$ 1 phút & ruff, pytest \\
    \bottomrule
  \end{tabularx}
\end{table}

\section{Definition of Done}
\label{sec:dod}

Một story chỉ được tính là Done khi thoả cả tám điều kiện dưới đây. Danh sách này áp dụng chung cho mọi story,
và là căn cứ để Product Owner chấp nhận story ở Sprint Review.
\begin{itemize}[label=$\square$]
  \item Code đã merge vào nhánh feature của sprint, có review của ít nhất 1 người khác.
  \item Đạt toàn bộ acceptance criteria; PO đã xem demo.
  \item Có test tự động cho logic mới (pytest hoặc Playwright); toàn bộ test cũ vẫn đạt.
  \item \texttt{ruff check} sạch; frontend \texttt{npm run build} (gồm type check) đạt.
  \item Không còn Bug mức Critical/High liên quan tới story.
  \item README/ARCHITECTURE được cập nhật nếu thay đổi cách chạy, cấu hình hoặc kiến trúc.
  \item Không có bí mật (key, mật khẩu) trong commit.
  \item Task trên Azure Boards đã cập nhật giờ thực tế và chuyển trạng thái Done.
\end{itemize}
Nhóm đã có DoD và các task QC riêng: \canDien[đối chiếu DoD ở trên với DoD đã thống nhất, giữ bản của nhóm nếu
khác].

\section{Rà soát chất lượng}

Ngoài kiểm thử, ba hoạt động rà soát gắn với quy trình Scrum giúp phát hiện lỗi sớm và tránh lặp lại lỗi cũ:
\begin{itemize}
  \item \textbf{Code review:} mọi thay đổi đi qua Pull Request trên nhánh feature trước khi merge.
  \item \textbf{Sprint Review:} demo trên bản build thật (bản web hoặc bộ cài), không demo trên máy dev đã sửa tay.
  \item \textbf{Retrospective:} xem lại các lỗi lọt qua test và bổ sung test tự động cho đúng lỗi đó.
\end{itemize}

% ============================================================ Chapter 3
\chapter{Kế hoạch kiểm thử}
\label{chap:testplan}

\textbf{Kiểm thử được tổ chức theo tám cấp, phần lớn tự động.} Unit test và test API chạy ở mỗi commit; test E2E
trên trình duyệt và desktop chạy mỗi sprint hoặc trước khi build bộ cài; test thủ công và UAT chốt chặn cuối trước
phát hành. Một bản chỉ được phát hành khi không còn lỗi Critical/High và PO, Sponsor đồng ý.

\section{Chiến lược kiểm thử}

Bảng~\ref{tab:strategy} mô tả phạm vi, công cụ, người thực hiện và thời điểm của từng cấp kiểm thử. Các phụ
thuộc bên ngoài (STT, LLM, S3) được thay bằng bản giả trong test tự động, nên bộ test chạy nhanh và cho kết quả
ổn định.

\begin{table}[htbp]
  \centering
  \caption{Chiến lược kiểm thử theo cấp}
  \label{tab:strategy}
  \small
  \setlength{\tabcolsep}{4pt}
  \begin{tabularx}{\linewidth}{@{}P{2.3cm}LP{3.5cm}P{1.9cm}P{2.2cm}@{}}
    \toprule
    \thead{Cấp} & \thead{Phạm vi} & \thead{Công cụ} & \thead{Ai} & \thead{Khi nào} \\
    \midrule
    Unit / component & Logic live (LocalAgreement), chia đoạn file dài, parse API STT, prompt/JSON LLM, xuất file
      & pytest (STT giả, LLM mock, S3 giả lập bằng moto) & Dev & Mỗi commit \\
    API / tích hợp & Endpoint meetings, uploads, summary, translation, \path{/v1}, S3 & pytest, FastAPI TestClient
      & Dev & Mỗi commit \\
    E2E trình duyệt & Ghi âm (micro, micro + âm thanh họp), offline rồi đồng bộ, upload, file sai định dạng; tóm
      tắt, dịch (Sprint 3) & Playwright trên Chrome và Edge thật, file WAV giả làm micro & QC & Mỗi sprint \\
    E2E desktop & Cửa sổ, micro, âm thanh hệ thống, thoát sạch, bản đóng gói
      & Playwright Electron (\path{desktop/scripts/smoke-test.js}) & Dev desktop & Trước khi build bộ cài \\
    Hiệu năng & Độ trễ live, thời gian xử lý file dài & Script mô phỏng ghi âm thời gian thực, upload file dài
      & Dev STT & Khi đổi engine hoặc model \\
    Chất lượng AI & ROUGE-L, độ phủ quyết định, chrF++, BLEU, giữ thuật ngữ & \path{backend/scripts/eval_*.py}
      & Dev AI, BA (mẫu tham chiếu) & Khi đổi model hoặc prompt \\
    Thủ công / UAT & Test case~\cite{testcases}, điện thoại, Firefox/Safari, bản web AWS & Thao tác tay
      & QC, người dùng thử & Cuối sprint, trước phát hành \\
    Deploy & Build image, HTTPS, đăng nhập, \path{/v1}, S3 & Docker local, checklist hướng dẫn deploy
      & Dev ops & Trước khi deploy \\
    \bottomrule
  \end{tabularx}
\end{table}

Hình~\ref{fig:test-layers} xếp các cấp trên thành ba tầng theo tần suất chạy. Tầng dưới cùng chạy nhiều nhất và
rẻ nhất; tầng trên cùng chạy ít nhất nhưng là nơi duy nhất có người dùng thật. Ba loại kiểm tra chuyên biệt
(hiệu năng, chất lượng AI, deploy) không chạy theo lịch cố định mà theo sự kiện.

\begin{figure}[htbp]
  \centering
  \begin{tikzpicture}[
      layer/.style={draw=primaryblue,fill=accentblue,align=center,font=\footnotesize,
                    minimum height=1.25cm,inner sep=3pt,execute at begin node=\setstretch{1.05}},
      note/.style={font=\footnotesize\itshape,text=chartgray,anchor=west,align=left,
                   execute at begin node=\setstretch{1.05}},
      lead/.style={chartgray,densely dotted}]
    % Three stacked layers, widest (most frequent, cheapest) at the bottom
    \node[layer,minimum width=10cm] (l1) at (5,0.625)
      {\textbf{Unit / component, API / tích hợp}\\ pytest với STT giả, LLM mock, S3 giả lập bằng moto};
    \node[layer,minimum width=7.6cm,fill=accentblue!55] (l2) at (5,1.875)
      {\textbf{E2E trình duyệt, E2E desktop}\\ Playwright trên Chrome, Edge và Electron};
    \node[layer,minimum width=5.2cm,draw=todored,fill=todored!10] (l3) at (5,3.125)
      {\textbf{Thủ công, UAT}\\ QC và người dùng thử};
    % When each layer runs
    \draw[lead] (l1.east) -- (10.6,0.625);
    \draw[lead] (l2.east) -- (10.6,1.875);
    \draw[lead] (l3.east) -- (10.6,3.125);
    \node[note] at (10.6,0.625) {Mỗi commit\\ (tự động)};
    \node[note] at (10.6,1.875) {Mỗi sprint, trước khi\\ build bộ cài (tự động)};
    \node[note] at (10.6,3.125) {Cuối sprint, trước phát\\ hành (thao tác tay)};
    % Event-driven checks
    \node[draw=chartgray,dashed,rounded corners=2pt,align=left,font=\footnotesize,text width=14cm,
          inner sep=4pt,anchor=north west] at (0,-0.35)
      {\textbf{Theo sự kiện:} hiệu năng khi đổi engine hoặc model; chất lượng AI khi đổi model hoặc prompt; kiểm
       tra deploy trước mỗi lần deploy.};
  \end{tikzpicture}
  \caption{Ba tầng kiểm thử theo tần suất chạy và các kiểm tra theo sự kiện}
  \label{fig:test-layers}
\end{figure}

\section{Quản lý lỗi}

Mỗi lỗi được xếp một trong bốn mức theo mức ảnh hưởng tới người dùng; hai mức cao nhất chặn phát hành
(Bảng~\ref{tab:severity}).

\begin{table}[htbp]
  \centering
  \caption{Mức độ lỗi và thời hạn xử lý}
  \label{tab:severity}
  \begin{tabularx}{\linewidth}{@{}P{2.2cm}LP{4.6cm}@{}}
    \toprule
    \thead{Mức} & \thead{Định nghĩa} & \thead{Thời hạn xử lý} \\
    \midrule
    Critical & Mất dữ liệu ghi âm, app không chạy, lộ dữ liệu & Ngay, chặn phát hành \\
    High & Tính năng chính hỏng, không có cách né & Trong sprint, chặn phát hành \\
    Medium & Lỗi có cách né, sai hiển thị quan trọng & Sprint sau \\
    Low & Thẩm mỹ, câu chữ & Khi có thời gian \\
    \bottomrule
  \end{tabularx}
\end{table}

Quy trình xử lý lỗi gồm bốn bước:
\begin{enumerate}
  \item QC tạo Bug trên Azure Boards, link tới story, kèm bước tái hiện, kết quả mong đợi, log hoặc ảnh chụp.
  \item Dev sửa lỗi.
  \item QC kiểm tra lại trên bản đã sửa.
  \item QC đóng Bug.
\end{enumerate}

\section{Tiêu chí bắt đầu, kết thúc và phát hành}
\label{sec:criteria}

Ba cửa kiểm soát quyết định khi nào bắt đầu test một sprint, khi nào dừng, và khi nào được phát hành
(Bảng~\ref{tab:criteria}). Tiêu chí phát hành được dùng để đánh giá bản v0.4 ở Mục~\ref{sec:release-eval}.

% Last table of the chapter: [H] so it is not pushed onto a page of its own.
\begin{table}[H]
  \centering
  \caption{Tiêu chí bắt đầu, kết thúc và phát hành}
  \label{tab:criteria}
  \begin{tabularx}{\linewidth}{@{}P{3.8cm}L@{}}
    \toprule
    \thead{Tiêu chí} & \thead{Nội dung} \\
    \midrule
    Bắt đầu test sprint & Story đạt DoD phía Dev, bản build chạy được, test case đã viết \\
    Kết thúc test sprint & 100\% test case ưu tiên cao đã chạy; tỷ lệ đạt $\geq$ 95\%; không còn lỗi Critical/High \\
    Phát hành (release) & Đạt tiêu chí kết thúc; có release notes~\cite{pcr}; smoke test bản đóng gói và bản web đạt;
      PO và Sponsor đồng ý \\
    \bottomrule
  \end{tabularx}
\end{table}

% ============================================================ Chapter 4
\chapter{Báo cáo tổng kết kiểm thử}
\label{chap:summary}

\textbf{Kết luận (12/11/2026): bản v0.4 đủ điều kiện phát hành cho đợt test nội bộ và UAT, nhưng chưa đủ để
nghiệm thu cuối.} Mọi test tự động đã chạy đều đạt và hiệu năng đạt mục tiêu; bộ cài desktop đã được build lại và
đạt smoke test. Còn thiếu ba việc: chạy lại test E2E trình duyệt sau thay đổi của Sprint 4, deploy thật lên AWS,
và đo độ chính xác tiếng Việt trên mẫu UAT.

\section{Phạm vi và môi trường}

Bảng~\ref{tab:environment} ghi phiên bản, máy và model được dùng khi kiểm thử. Toàn bộ số liệu trong chương được
đo trên cấu hình này.

\begin{table}[htbp]
  \centering
  \caption{Phạm vi và môi trường kiểm thử}
  \label{tab:environment}
  \begin{tabularx}{\linewidth}{@{}P{3.8cm}L@{}}
    \toprule
    \thead{Mục} & \thead{Nội dung} \\
    \midrule
    Phiên bản & Nhánh \path{docs/project-management} (gồm toàn bộ code tới \path{feature/aws-deploy}) \\
    Máy test & Windows 11, i7-13620H, RTX 4050 6\,GB, Python 3.13, Node 22, Chrome/Edge, Electron 45 \\
    Model & Whisper \texttt{large-v3-turbo} (GPU), \texttt{small} (CPU); LLM mock trong test tự động \\
    Kế hoạch tham chiếu & Chương~\ref{chap:testplan} và test case Sprint 3 (Phụ lục~\ref{app:testcases}) \\
    \bottomrule
  \end{tabularx}
\end{table}

\section{Kết quả test tự động}

Kết quả của tám bộ test tự động được tổng hợp trong Bảng~\ref{tab:auto-results}. Sáu bộ đã có kết quả đều đạt
100\%; hai bộ Playwright trên trình duyệt còn chờ chạy lại.

\begin{table}[htbp]
  \centering
  \caption{Kết quả test tự động trên bản v0.4}
  \label{tab:auto-results}
  \small
  \setlength{\tabcolsep}{4pt}
  \begin{tabularx}{\linewidth}{@{}LQ{1.5cm}Q{1.5cm}Q{1.6cm}P{4.4cm}@{}}
    \toprule
    \thead{Bộ test} & \thead{Số test} & \thead{Đạt} & \thead{Không đạt} & \thead{Ghi chú} \\
    \midrule
    pytest backend (API, live, STT remote, \path{/v1}, S3, tóm tắt, dịch) & 88 & 88 & 0 & Khoảng 6\,s \\
    ruff (lint) & -- & Sạch & -- & \\
    Type check và build frontend & -- & Đạt & -- & \\
    Playwright desktop -- từ mã nguồn & 8 & 8 & 0 & Cửa sổ, micro, âm thanh hệ thống, thoát sạch \\
    Playwright desktop -- bản đóng gói & 8 & 8 & 0 & Bộ cài v0.4 build lại ngày 12/11/2026 (205\,MB), có
      streaming và STT remote \\
    Docker: build image, HTTPS, đăng nhập, \path{/v1} & 8 kiểm tra & 8 & 0 & Stack chạy local với
      \texttt{SITE\_ADDRESS=}\allowbreak\texttt{localhost} \\
    \midrule
    Playwright trình duyệt (Chrome, Edge) -- E1 & 5 $\times$ 2 & \canDien & & Cần chạy lại sau khi đổi chunk 1\,s \\
    Playwright Sprint 3 (\path{sprint3.spec.ts}) & \canDien & & & Server \texttt{AI\_MOCK=true} \\
    \bottomrule
  \end{tabularx}
\end{table}

Hình~\ref{fig:auto-results} minh hoạ cùng số liệu theo tỷ lệ. Bộ Playwright Sprint 3 chưa có số test nên không
được vẽ.

\begin{figure}[htbp]
  \centering
  \begin{tikzpicture}
    \begin{axis}[
        width=0.78\linewidth, height=5.4cm,
        xbar stacked, bar width=0.42cm,
        xmin=0, xmax=100, xtick={0,25,50,75,100},
        xticklabel={\pgfmathprintnumber{\tick}\%},
        symbolic y coords={PY,DS,DP,DK,E1}, ytick=data, y dir=reverse,
        yticklabels={{pytest backend (88)},{Desktop -- mã nguồn (8)},{Desktop -- đóng gói (8)},
                     {Docker (8)},{Trình duyệt E1 (10)}},
        legend style={at={(0.5,-0.24)},anchor=north,legend columns=2,font=\small,draw=gray!40,
                      /tikz/every even column/.append style={column sep=0.6cm}},
        tick label style={font=\small}, enlarge y limits=0.15,
        grid=major, ymajorgrids=false, grid style={gray!25}]
      \addplot[fill=primaryblue, draw=primaryblue] coordinates {(100,PY) (100,DS) (100,DP) (100,DK) (0,E1)};
      \addlegendentry{Đạt}
      \addplot[fill=gray!25, draw=chartgray] coordinates {(0,PY) (0,DS) (0,DP) (0,DK) (100,E1)};
      \addlegendentry{Chưa có kết quả (cần chạy lại)}
    \end{axis}
  \end{tikzpicture}
  \caption{Tỷ lệ đạt của các bộ test tự động (số trong ngoặc là số test)}
  \label{fig:auto-results}
\end{figure}

\section{Kết quả hiệu năng}

Năm kịch bản hiệu năng được đo thực tế (Bảng~\ref{tab:perf}); tất cả đều đạt. Độ trễ live như nhau khi chạy GPU
trên máy và khi dùng engine remote, nhưng với engine remote máy khách không phải tải model.

\begin{table}[htbp]
  \centering
  \caption{Kết quả hiệu năng đo thực tế}
  \label{tab:perf}
  \small
  \begin{tabularx}{\linewidth}{@{}P{4.4cm}P{2.5cm}LQ{0.9cm}@{}}
    \toprule
    \thead{Kịch bản} & \thead{Mục tiêu} & \thead{Kết quả} & \thead{Đạt} \\
    \midrule
    Live, GPU local, file 17\,s phát theo thời gian thực & $\leq$ 3--5\,s
      & Từ mới khoảng 1\,s; câu chốt trung bình 1,93\,s, tối đa 2,31\,s & \pass \\
    Live qua engine remote (máy không có model) & Như trên
      & Trung bình 1,93\,s, tối đa 2,31\,s; máy khách không tải model & \pass \\
    File 11,6 phút -- remote & 1 giờ $\leq$ 10 phút & 17,8\,s (39$\times$ realtime), dòng đầu sau 3,1\,s & \pass \\
    File 11,6 phút -- GPU local & Như trên & 56,7\,s (12$\times$ realtime), dòng đầu sau 4,6\,s & \pass \\
    Khởi động app desktop & $\leq$ 10\,s & 1,4--1,9\,s & \pass \\
    \bottomrule
  \end{tabularx}
\end{table}

Để so sánh các kịch bản có đơn vị khác nhau, Hình~\ref{fig:perf-ratio} quy mỗi kết quả \emph{xấu nhất} về phần
trăm của ngưỡng mục tiêu, với ngưỡng live lấy cận chặt 3\,s. Mục tiêu ``1 giờ audio $\leq$ 10 phút'' tương ứng
với ngưỡng $11{,}6 / 6 \approx 1{,}93$ phút (khoảng 116\,s) cho file 11,6 phút. Mọi cột đều dưới 100\%; độ trễ live
là chỉ số sát ngưỡng nhất, ở mức khoảng 77\%.

\begin{figure}[htbp]
  \centering
  \begin{tikzpicture}
    \begin{axis}[
        width=0.74\linewidth, height=5.6cm,
        xbar, bar width=0.42cm,
        xmin=0, xmax=110, xtick={0,25,50,75,100},
        xticklabel={\pgfmathprintnumber{\tick}\%},
        symbolic y coords={LL,LR,FR,FL,AP}, ytick=data, y dir=reverse,
        yticklabels={{Live, GPU local},{Live, engine remote},{File 11,6 phút, remote},
                     {File 11,6 phút, GPU local},{Khởi động app desktop}},
        xlabel={Kết quả xấu nhất / ngưỡng mục tiêu},
        nodes near coords, nodes near coords={\pgfmathprintnumber[fixed,precision=0]{\pgfplotspointmeta}\%},
        every node near coord/.append style={font=\footnotesize},
        tick label style={font=\small}, enlarge y limits=0.15,
        grid=major, ymajorgrids=false, grid style={gray!25},
        /pgf/number format/.cd, use comma, 1000 sep={.}]
      \addplot[fill=primaryblue, draw=primaryblue]
        coordinates {(77,LL) (77,LR) (15.3,FR) (48.9,FL) (19,AP)};
      % x range is 0..110, so 100% sits at 100/110 of the axis width.
      \draw[todored, dashed, thick] (rel axis cs:0.90909,0) -- (rel axis cs:0.90909,1);
    \end{axis}
  \end{tikzpicture}
  \caption{Kết quả hiệu năng theo phần trăm ngưỡng mục tiêu (đường đỏ: 100\%)}
  \label{fig:perf-ratio}
\end{figure}

\section{Kết quả test thủ công}

Kết quả test thủ công của QC được ghi trong Bảng~\ref{tab:manual}. Tại thời điểm báo cáo, các bộ test case thủ
công chưa có kết quả trong tài liệu; bảng tính kết quả kiểm thử của QC~\cite{tsxlsx} cũng đang ghi các mục kiểm tra
tay là ``Đang chờ''.

\begin{table}[htbp]
  \centering
  \caption{Kết quả test thủ công của QC}
  \label{tab:manual}
  \begin{tabularx}{\linewidth}{@{}LQ{2.3cm}Q{2.3cm}Q{1.8cm}Q{1.8cm}@{}}
    \toprule
    \thead{Bộ test case} & \thead{Tổng} & \thead{Đạt} & \thead{Không đạt} & \thead{Chưa chạy} \\
    \midrule
    \#99 Tóm tắt (S-01\ldots S-11) & 11 & \canDien & & \\
    \#100 Dịch (T-01\ldots T-12) & 12 & \canDien & & \\
    Điện thoại Android/iOS (T1.2.5) & \canDien & & & \\
    UAT bản web trên AWS & \canDien & & & \\
    \bottomrule
  \end{tabularx}
\end{table}

\section{Lỗi còn tồn và hạn chế đã biết}

Bốn hạn chế đã biết đều ở mức Medium hoặc Low, nên không chặn phát hành (Bảng~\ref{tab:defects}). Hai hạn chế
đầu tương ứng với rủi ro R13 và R12 đã được chấp nhận trong Risk Register~\cite{pmp}.

\begin{table}[htbp]
  \centering
  \caption{Lỗi còn tồn và hạn chế đã biết}
  \label{tab:defects}
  \small
  \begin{tabularx}{\linewidth}{@{}P{1.3cm}LQ{1.6cm}P{4.4cm}@{}}
    \toprule
    \thead{ID} & \thead{Mô tả} & \thead{Mức} & \thead{Trạng thái} \\
    \midrule
    L-01 & iOS dừng micro khi khoá màn hình (giới hạn của trình duyệt) & Medium & Chấp nhận (R13) \\
    L-02 & Bộ cài chưa ký số, SmartScreen cảnh báo & Low & Chấp nhận (R12) \\
    L-03 & Engine remote dùng API công cộng: ngôn ngữ được nhận theo block khoảng 30\,s, thô hơn engine local (theo câu)
      & Low & Dùng server SumTalk GPU để giữ nhận ngôn ngữ theo câu \\
    L-04 & Live: dòng có thể bị ngắt giữa câu khi nói liền $>$ 8\,s không nghỉ (bản cuối sau khi Stop sẽ chia lại câu)
      & Low & Thiết kế có chủ ý \\
    \canDien & Lỗi QC tìm thấy trên Azure Boards & & \\
    \bottomrule
  \end{tabularx}
\end{table}

\section{Đánh giá theo tiêu chí phát hành}
\label{sec:release-eval}

Bảng~\ref{tab:release-eval} đối chiếu bản v0.4 với tiêu chí phát hành ở Mục~\ref{sec:criteria}. Chỉ tiêu chí hiệu năng
đạt trọn vẹn; ba tiêu chí đạt một phần hoặc còn chờ xác nhận, và chưa có sự đồng ý của PO, Sponsor.

\begin{table}[htbp]
  \centering
  \caption{Đánh giá bản v0.4 theo tiêu chí phát hành}
  \label{tab:release-eval}
  \begin{tabularx}{\linewidth}{@{}P{5.4cm}L@{}}
    \toprule
    \thead{Tiêu chí} & \thead{Trạng thái} \\
    \midrule
    Test tự động đạt 100\% & \pass{} Backend, desktop, Docker; E2E trình duyệt cần chạy lại \\
    Không còn lỗi Critical/High & \pass{} Theo hiểu biết hiện tại; \canDien[xác nhận trên board] \\
    Hiệu năng đạt mục tiêu & \pass \\
    Smoke test bản đóng gói và bản web & Bản đóng gói: \pass{} (8/8); bản web: đạt khi chạy local, chưa có trên AWS \\
    PO và Sponsor đồng ý & \canDien \\
    \bottomrule
  \end{tabularx}
\end{table}

\heading{Khuyến nghị}
Chạy lại test E2E trình duyệt và bộ test Sprint 3, deploy lên AWS, chạy UAT trong 1 tuần, rồi nghiệm thu bản v1.0.

% ============================================================ Chapter 5
\chapter{Kết luận và phê duyệt}
\label{chap:conclusion}

\section{Tóm tắt}

Tài liệu đã đặt kế hoạch chất lượng và kiểm thử cạnh kết quả trên bản v0.4:
\begin{itemize}
  \item \textbf{Kế hoạch:} năm mục tiêu chất lượng, Definition of Done gồm tám điều kiện, tám cấp kiểm thử chia
        thành ba tầng theo tần suất, bốn mức lỗi và ba cửa kiểm soát bắt đầu, kết thúc, phát hành.
  \item \textbf{Kết quả tự động:} 88/88 test backend, 8/8 smoke test desktop (mã nguồn và bản đóng gói), 8/8 kiểm
        tra Docker, lint và build frontend đạt.
  \item \textbf{Hiệu năng:} cả năm kịch bản đạt; kết quả xấu nhất của độ trễ live là 2,31\,s, khoảng 77\% ngưỡng 3\,s.
  \item \textbf{Kết luận:} v0.4 đủ điều kiện cho test nội bộ và UAT; chưa đủ để nghiệm thu cuối.
\end{itemize}

\section{Việc cần làm trước khi nghiệm thu}

\begin{enumerate}
  \item Chạy lại Playwright trình duyệt (E1, Chrome và Edge) và bộ \path{sprint3.spec.ts}, rồi điền
        Bảng~\ref{tab:auto-results}.
  \item QC chạy 23 test case S-01\ldots S-11, T-01\ldots T-12 và test trên điện thoại; điền Bảng~\ref{tab:manual} và
        bảng kết quả ở Phụ lục~\ref{app:testcases}.
  \item Deploy bản web lên AWS, chạy smoke test và UAT trong 1 tuần; đo độ chính xác tiếng Việt trên mẫu UAT.
  \item Xác nhận trên Azure Boards không còn lỗi Critical/High; bổ sung lỗi QC tìm thấy vào Bảng~\ref{tab:defects}.
  \item Đối chiếu Definition of Done ở Mục~\ref{sec:dod} với DoD nhóm đã thống nhất; điền các ô \canDien{} còn lại.
\end{enumerate}

\needspace{14\baselineskip}
\section{Phê duyệt}

Nhóm dự án đề nghị QC, Product Owner và Project Manager phê duyệt kế hoạch chất lượng, kế hoạch kiểm thử và kết luận
kiểm thử của bản v0.4 trình bày trong tài liệu này.

\begin{table}[H]
  \centering
  \renewcommand{\arraystretch}{2.4}
  \begin{tabularx}{\linewidth}{@{}P{3.6cm}LP{3.2cm}P{2.6cm}@{}}
    \toprule
    \thead{Vai trò} & \thead{Họ tên} & \thead{Chữ ký} & \thead{Ngày} \\
    \midrule
    QC & & & \\
    Product Owner & & & \\
    Project Manager & & & \\
    \bottomrule
  \end{tabularx}
\end{table}

% ============================================================ References
\begin{thebibliography}{9}
\addcontentsline{toc}{chapter}{\bibname}
  \bibitem{srs} Nhóm SumTalk, \emph{Đặc tả yêu cầu phần mềm (Software Requirements Specification)},
    Tài liệu 02, mã SUMTALK-SRS, phiên bản 2.1, 24/11/2026.
  \bibitem{pmp} Nhóm SumTalk, \emph{Kế hoạch quản lý dự án (Project Management Plan)}, Tài liệu 03,
    mã SUMTALK-PMP, phiên bản 2.1, 24/11/2026.
  \bibitem{pcr} Nhóm SumTalk, \emph{Báo cáo kết thúc dự án (Project Closure Report)}, Tài liệu 06,
    mã SUMTALK-PCR, phiên bản 2.0, 13/11/2026.
  \bibitem{testcases} Nhóm SumTalk, \emph{Sprint 3 test cases (\#99, \#100)}, kho mã nguồn
    \texttt{DatTapTanhCode/SumTalk}, tệp \path{docs/testing/sprint3-test-cases.md}.
  \bibitem{tsxlsx} Nhóm SumTalk, \emph{Báo cáo kiểm thử SumTalk -- phiên bản v0.4} (bảng tính), kho mã nguồn
    \texttt{DatTapTanhCode/SumTalk}, tệp \path{TEST_SUMMARY_13112026_SumTalk_v2.xlsx}, 13/11/2026.
  \bibitem{rouge} C.-Y. Lin, ``ROUGE: A package for automatic evaluation of summaries,'' in \emph{Text
    Summarization Branches Out}, Barcelona: Association for Computational Linguistics, 2004, pp.~74--81.
  \bibitem{bleu} K. Papineni, S. Roukos, T. Ward, W.-J. Zhu, ``BLEU: a method for automatic evaluation of machine
    translation,'' in \emph{Proceedings of the 40th Annual Meeting of the ACL}, 2002, pp.~311--318.
  \bibitem{chrf} M. Popović, ``chrF++: words helping character n-grams,'' in \emph{Proceedings of the Second
    Conference on Machine Translation}, 2017, pp.~612--618.
\end{thebibliography}

% ============================================================ Appendix
\appendix
\chapter{Test case E2E Sprint 3 (\#99, \#100)}
\label{app:testcases}

Phụ lục này chép lại bộ test case QC dùng hằng ngày~\cite{testcases}; nội dung giữ nguyên tiếng Anh như bản QC
đang dùng. Mỗi test case ứng với một acceptance criterion (cột AC) của story \#67 hoặc \#60. Lỗi tìm thấy được ghi
thành \textbf{Bug} trên Azure Boards, link tới story và test lại sau khi sửa. Các chỉ số ROUGE-L~\cite{rouge},
BLEU~\cite{bleu} và chrF++~\cite{chrf} dùng trong \path{eval_summary.py} và \path{eval_translation.py} hỗ trợ cho
các test case chất lượng S-04 và T-04.

\section{Setup}

The three ways to run the tests are listed in Table~\ref{tab:setup}.

\begin{table}[htbp]
  \centering
  \caption{Test setup}
  \label{tab:setup}
  \small
  \begin{tabularx}{\linewidth}{@{}P{3.6cm}L@{}}
    \toprule
    \thead{Run} & \thead{How} \\
    \midrule
    \textbf{Mock model} (predictable, no GPU) & \texttt{AI\_MOCK=true} in \path{backend/.env}, then \path{start.bat}.
      Summary text starts with ``Bản tóm tắt thử nghiệm (mock)'' (vi) or ``Mock summary'' (en); a translated line is
      the original prefixed with \texttt{[vi]} / \texttt{[en]}. \\
    \textbf{Real model} (quality) & \texttt{AI\_MOCK=false}, \texttt{LLM\_BASE\_URL} + \texttt{SUMMARY\_MODEL} /
      \texttt{TRANSLATION\_MODEL} set, Ollama running (or cloud key). \\
    \textbf{Simulate a model failure} & \texttt{AI\_MOCK=false} and \path{LLM_BASE_URL=http://localhost:9/v1}
      (nothing listens there). \\
    \bottomrule
  \end{tabularx}
\end{table}

\textbf{Sample audio:} \path{frontend/e2e/fixtures/speech_en.mp3} (English, 17\,s), a Vietnamese recording, a mixed
EN/VI recording, and a recording of 60 minutes or more (any meeting recording; the content does not matter).

\textbf{Automated coverage:}
\begin{itemize}
  \item pytest, mock model: \path{backend/tests/test_summary_translation_api.py}, \path{test_ai_summary.py},
        \path{test_ai_translation.py};
  \item Playwright, server with \texttt{AI\_MOCK=true}: \path{frontend/e2e/sprint3.spec.ts}, run with
        \texttt{npx playwright test e2e/sprint3.spec.ts}.
\end{itemize}

\section{\#99 -- Summary (story \#67)}

\begingroup\small
\setlength{\tabcolsep}{4pt}
\begin{longtable}{@{}P{1.1cm}Q{0.7cm}P{4.5cm}P{5.8cm}P{2.5cm}@{}}
  \caption{Test cases \#99 -- Summary}\label{tab:tc-summary}\\
  \toprule
  \thead{ID} & \thead{AC} & \thead{Steps} & \thead{Expected} & \thead{Auto} \\
  \midrule
  \endfirsthead
  \multicolumn{5}{@{}l}{\footnotesize\itshape Bảng~\thetable{} (tiếp theo)}\\
  \toprule
  \thead{ID} & \thead{AC} & \thead{Steps} & \thead{Expected} & \thead{Auto} \\
  \midrule
  \endhead
  \bottomrule
  \endlastfoot
  S-01 & 1 & Upload \path{speech_en.mp3}. Wait for \ui{Ready}. Open the \ui{Summary} tab without clicking anything
    else. & ``Summarizing…'' with a spinner, then the summary appears by itself (page polls every 3\,s).
    & e2e, pytest \\
  S-02 & 2 & Look at a finished summary. & Sections \ui{Overview, Key points, Decisions, Action items, Open
    questions}. Empty sections say so (``No decisions were recorded.''). & e2e, pytest \\
  S-03 & 3 & Regenerate ▾ $\rightarrow$ \ui{English}, confirm. Then Regenerate ▾ $\rightarrow$ \ui{Vietnamese}.
    Repeat on the mixed EN/VI meeting. & Text is in the chosen language each time, also for the mixed meeting.
    English technical terms (API, deploy, sprint…) stay in English in the Vietnamese summary.
    & e2e (lang), manual (terms) \\
  S-04 & 3 & Real model: compare the summary with the transcript. & No people, numbers, dates or decisions that are
    not in the transcript. & manual / \path{eval_summary.py} \\
  S-05 & 4 & Upload a recording of \textbf{$\geq$ 60 min}. & Summary ends as \textbf{done} (no error). In
    \path{GET /api/meetings/{id}/summary}: \texttt{mode = "map\_reduce"}, \texttt{calls} = number of parts + 1.
    & pytest (80k-char transcript) \\
  S-06 & 5 & On a finished summary: Regenerate ▾ $\rightarrow$ a language. & A confirmation ``The current summary
    will be replaced…'' appears. After confirming, the old text disappears and the new one replaces it. Cancel
    keeps the old one. & e2e \\
  S-07 & 5 & While a summary is being generated. & The Regenerate button is disabled.
    \texttt{POST \ldots/summary/regenerate} returns \textbf{409}. & pytest \\
  S-08 & 6 & Simulate a model failure, upload a file. & After the retries (about 30\,s) the tab shows ``The summary
    could not be generated'' with the error and \ui{Try again}. Nothing spins forever. & pytest \\
  S-09 & 6 & Fix the model URL, click \ui{Try again}. & The summary is generated. The transcript is \textbf{not}
    re-transcribed (meeting stays Ready). & pytest \\
  S-10 & -- & Re-transcribe a meeting that has a summary. & The summary is rebuilt after the new transcript.
    & pytest \\
  S-11 & -- & Phone width ($\leq$ 400\,px). & Sections stack in one column, buttons reachable, no horizontal
    scroll. & manual \\
\end{longtable}
\endgroup

\section{\#100 -- Translation (story \#60)}

\begingroup\small
\setlength{\tabcolsep}{4pt}
\begin{longtable}{@{}P{1.1cm}Q{0.7cm}P{4.5cm}P{5.8cm}P{2.5cm}@{}}
  \caption{Test cases \#100 -- Translation}\label{tab:tc-translation}\\
  \toprule
  \thead{ID} & \thead{AC} & \thead{Steps} & \thead{Expected} & \thead{Auto} \\
  \midrule
  \endfirsthead
  \multicolumn{5}{@{}l}{\footnotesize\itshape Bảng~\thetable{} (tiếp theo)}\\
  \toprule
  \thead{ID} & \thead{AC} & \thead{Steps} & \thead{Expected} & \thead{Auto} \\
  \midrule
  \endhead
  \bottomrule
  \endlastfoot
  T-01 & 1 & Open the \ui{Translation} tab of a Ready meeting. Open ``Translate to''. & The list offers at least
    \ui{Vietnamese} and \ui{English}; choosing one starts the translation. & e2e \\
  T-02 & 2 & Translate \path{speech_en.mp3} into Vietnamese; wait for the end. & A progress bar with \% while
    running, lines filling in. When done, \textbf{every} line has a translation (\texttt{GET \ldots/translations/vi}: no
    \texttt{translation: null}). & e2e, pytest \\
  T-03 & 3 & Translate the English meeting into \ui{English}, and the Vietnamese meeting into \ui{Vietnamese}.
    & Lines are kept unchanged with the label ``already English/Vietnamese'' (\texttt{skipped: true}), and no model
    call is made for them. & pytest \\
  T-04 & 3 & Real model: translate the mixed EN/VI meeting into Vietnamese and into English. & Natural translation;
    English technical terms and names kept. Vietnamese lines are kept when the target is Vietnamese.
    & manual / \path{eval_translation.py} \\
  T-05 & 4 & Look at a finished translation on a desktop screen. & Two columns \ui{Original \textbar{}
    Translation}, one row per segment, time on the left, rows aligned. Clicking a time plays the audio from there.
    & e2e \\
  T-06 & 4 & Same at phone width. & Each original line is directly above its translation. & manual \\
  T-07 & 5 & Reload the page (F5), open the Translation tab. & The saved translation shows at once. No new job:
    \texttt{GET \ldots/translations} still has one entry with the same \texttt{updated\_at}; choosing the same language
    again returns \textbf{200} immediately. & e2e, pytest \\
  T-08 & 6 & Export ▾ $\rightarrow$ \ui{Text (.txt)}; open in Windows Notepad. & File \texttt{<title>\_vi.txt};
    Vietnamese accents display correctly; lines \texttt{[hh:mm:ss] translation}. & e2e (file), manual (Notepad) \\
  T-09 & 6 & Export ▾ $\rightarrow$ \ui{Subtitles (.srt)}; open in VLC with the audio, or in a subtitle editor.
    & File \texttt{<title>\_vi.srt}; numbered cues \texttt{HH:MM:SS,mmm -{}-> HH:MM:SS,mmm}; subtitles in sync.
    & e2e, pytest \\
  T-10 & 6 & Export the bilingual variants. & Original line above each translation, in TXT and SRT. & pytest \\
  T-11 & -- & Simulate a model failure, pick a language. & ``The translation stopped'' with the error and
    \ui{Try again}; Export disabled. After fixing the model, Try again translates only the missing lines. & pytest \\
  T-12 & -- & Before the transcript is Ready. & The tab says translation is available once the transcript is
    ready; the API returns 409. & pytest \\
\end{longtable}
\endgroup

\section{Results}

QC records each test run in Table~\ref{tab:tc-results}.

\begin{table}[H]
  \centering
  \caption{Test run results}
  \label{tab:tc-results}
  \small
  \renewcommand{\arraystretch}{1.8}
  \begin{tabularx}{\linewidth}{@{}P{1.8cm}P{2.6cm}P{2.2cm}P{2.2cm}Q{1.5cm}L@{}}
    \toprule
    \thead{Date} & \thead{Build / commit} & \thead{Tester} & \thead{Model} & \thead{Passed} & \thead{Failed (bug IDs)} \\
    \midrule
    & & & & & \\
    & & & & & \\
    & & & & & \\
    \bottomrule
  \end{tabularx}
\end{table}

\end{document}
